\chapter[\emj{🤑}~Greedy (Algoritmos Golosos)]{\emj{🤑}~Greedy (Algoritmos Golosos)}

\begin{dsaquees}

Tomar SIEMPRE la mejor opción del momento 🍰, sin pensar en el futuro,
esperando que esas decisiones locales sumen la mejor solución global.

Si quieres dar cambio con la menor cantidad de monedas, tomas la moneda
más grande que quepa una y otra vez. Eso es ser "goloso": eliges lo mejor
AHORA y no reconsideras. A veces funciona perfecto (monedas normales),
a veces falla (monedas raras), y ahí está el arte: saber CUÁNDO un
problema se puede resolver golosamente.

\end{dsaquees}

\begin{dsanino}

Imagina que tienes una mochila y quieres meter la mayor cantidad de
dulces 🍬, pero solo puedes agarrar de a uno rapidísimo. La regla golosa
dice: "agarra el dulce más grande que veas, no lo pienses". Sigues
agarrando el más grande que quede, hasta llenarte.

A veces esa regla te da lo mejor. Pero ojo: si los dulces vinieran
pegados de formas raras, agarrar siempre el más grande podría dejarte sin
espacio para dos medianos que juntos valían más. Por eso, antes de ser
goloso, hay que estar SEGURO de que "lo mejor ahorita" también es "lo
mejor al final". Cuando eso es cierto, ser goloso es rapidísimo y fácil.

\end{dsanino}

\begin{dsaotraforma}

Un algoritmo greedy construye la solución paso a paso, eligiendo en cada
etapa la opción que se ve mejor localmente (según un criterio), sin
retroceder.

Solo da el óptimo si el problema cumple dos propiedades:
\begin{itemize}
\item \textbf{Greedy choice property}: una elección local óptima lleva a
      una solución global óptima.
\item \textbf{Subestructura óptima}: la solución óptima contiene
      soluciones óptimas de sus subproblemas.
\end{itemize}

Si NO se cumplen, greedy falla y hay que usar Programación Dinámica.
Ejemplos donde greedy SÍ funciona: activity selection, intervalos,
código de Huffman, Dijkstra, MST (Kruskal/Prim), fractional knapsack.

\end{dsaotraforma}

\begin{dsadiagrama}
ACTIVITY SELECTION (máximo de reuniones sin traslape):

Reuniones (inicio, fin):
  A(1,3)  B(2,5)  C(4,7)  D(6,8)  E(5,9)

Paso 1: ordena por HORA DE FIN:  A(3) B(5) C(7) D(8) E(9)
Paso 2: toma A. Última fin = 3.
        B empieza en 2 < 3 -> choca, salta.
        C empieza en 4 >= 3 -> toma C. Última fin = 7.
        D empieza en 6 < 7 -> choca, salta.
        E empieza en 5 < 7 -> choca, salta.
Resultado: {A, C} -> 2 reuniones (óptimo).

Criterio goloso ganador: SIEMPRE la que TERMINA más temprano.
\end{dsadiagrama}

\begin{lstlisting}
PATRÓN GREEDY (3 pasos)
1. Encuentra el criterio de ordenamiento correcto (el "más pequeño",
   el que "termina antes", el "más valioso por kilo"...).
2. Ordena por ese criterio.
3. Recorre una vez tomando lo que quepa/no choque.

PSEUDOCÓDIGO (activity selection)
──────────────────────────────────────────────────────────────

función maxActividades(intervalos):
    ordenar intervalos por hora de FIN ascendente
    finAnterior = -infinito; conteo = 0
    PARA cada (inicio, fin) en intervalos:
        SI inicio >= finAnterior:
            conteo++; finAnterior = fin       // se toma
    return conteo

CÓDIGO C++ (activity selection / non-overlapping intervals)
──────────────────────────────────────────────────────────────

#include <vector>
#include <algorithm>
using namespace std;

int maxActividades(vector<pair<int,int>>& iv) {   // {inicio, fin}
    sort(iv.begin(), iv.end(),
         [](auto& a, auto& b){ return a.second < b.second; }); // por fin
    int fin = INT_MIN, cont = 0;
    for (auto& [ini, f] : iv) {
        if (ini >= fin) { cont++; fin = f; }      // no choca -> tómala
    }
    return cont;
}

CÓDIGO C++ (fractional knapsack: greedy SÍ funciona)
──────────────────────────────────────────────────────────────

// A diferencia del 0/1 knapsack (que es DP), aquí puedes partir objetos.
double fractionalKnapsack(int W, vector<pair<int,int>>& items) {
    // items = {valor, peso}; ordena por valor/peso DESC
    sort(items.begin(), items.end(), [](auto& a, auto& b){
        return (double)a.first/a.second > (double)b.first/b.second;
    });
    double total = 0;
    for (auto& [val, peso] : items) {
        if (W >= peso) { total += val; W -= peso; }   // cabe entero
        else { total += val * ((double)W / peso); break; } // fracción
    }
    return total;
}

COMPLEJIDAD (Big-O)
──────────────────────────────────────────────────────────────

  Problema            │ Tiempo      │ Domina
  ────────────────────┼─────────────┼──────────────
  Activity selection  │ O(n log n)  │ el sort
  Fractional knapsack │ O(n log n)  │ el sort
  Recorrido greedy    │ O(n)        │ una pasada

* Casi siempre el costo es el ORDENAMIENTO inicial.
\end{lstlisting}

\begin{lstlisting}
💡 Problemas greedy clásicos y su criterio ganador:
   - Intervalos sin traslape (LC 435): ordena por FIN, quita los que chocan.
   - Merge Intervals (LC 56): ordena por INICIO, fusiona solapados.
   - Jump Game (LC 55): mantén el alcance máximo; si lo superas, pierdes.
   - Gas Station (LC 134): si el tanque total >= 0, hay solución única.
   - Huffman: combina siempre los dos nodos de menor frecuencia (min-heap).
\end{lstlisting}

\begin{dsaentrevistas}

✅ El reto es DEMOSTRAR que greedy funciona. Si no puedes justificar por
   qué la elección local da el óptimo global, probablemente NO sea greedy
   y necesites DP. No lo uses "de fe".

💡 Contraejemplo de monedas: con monedas {1, 3, 4} y cambio de 6, greedy
   da 4+1+1 (3 monedas) pero lo óptimo es 3+3 (2 monedas). Por eso "coin
   change" general es DP, no greedy.

⚠️ El criterio de ordenamiento lo es TODO. En activity selection ordenar
   por duración o por inicio FALLA; solo ordenar por hora de fin da el
   óptimo. Piensa bien el criterio.

🔑 Greedy vs DP: si "elegir lo mejor ahora" nunca te arrepiente ->
   greedy (rápido). Si necesitas explorar combinaciones y recordar
   subproblemas -> DP.

\end{dsaentrevistas}

\begin{dsaresumen}

• Greedy: elige lo mejor local en cada paso, sin retroceder.
• Solo da el óptimo con greedy-choice + subestructura óptima.
• Patrón: encuentra criterio -> ordena -> recorre tomando lo que quepa.
• Activity selection: ordena por hora de FIN.
• Fractional knapsack SÍ es greedy; el 0/1 knapsack NO (es DP).
• Si no puedes demostrar que funciona, usa DP.
════════════════════════════════════════════════════════════════

\end{dsaresumen}
